\section{Acronym Generation}\label{section:acronym}

Good acronyms provide a concise and memorable way to communicate complex ideas, making them easier to understand and remember, ultimately leading to more efficient and effective communication.
Like in email writing, acronym generation also requires an iterative refinement process to achieve a concise and memorable representation of a complex term or phrase. 
Acronyms often involve tradeoffs between length, ease of pronunciation, and relevance to the original term or phrase.
Thus, acronym generation is a natural method testbed for our approach.

We source the dataset for this task from \url{https://github.com/krishnakt031990/Crawl-Wiki-For-Acronyms/blob/master/AcronymsFile.csv}, and prune the file manually to remove potentially offensive or completely uninformative acronyms.
This exercise generated a list of 250 acronyms.
The complete list is given in our code repository.


\paragraph{\fbmod} For feedback, we design an \fbmod that can provide multifaceted feedback. Specifically, each acronym is judged along five dimensions:
\begin{itemize}
\item \textbf{Ease of pronunciation:} How easy or difficult is it to pronounce the acronym? Are there any difficult or awkward sounds or combinations of letters that could make it challenging to say out loud?
\item \textbf{Ease of spelling:} How easy or difficult is it to spell the acronym? Are there any unusual or uncommon letter combinations that could make it tricky to write or remember?
\item \textbf{Relation to title:} How closely does the acronym reflect the content or topic of the associated title, phrase, or concept? Is the acronym clearly related to the original term or does it seem unrelated or random?
\item \textbf{Positive connotation:} Does the acronym have any positive or negative associations or connotations? Does it sound upbeat, neutral, or negative in tone or meaning?
\item \textbf{Well-known:} How familiar or recognizable is the acronym to the target audience? Is it a common or widely-used term, or is it obscure or unfamiliar?
\end{itemize}

Some of these criteria are difficult to quantify, and are a matter of human preference.
As with other modules, we leverage the superior instruction following capabilities of modern \llms to instead provide a few demonstrations of each task.
Crucially, the feedback includes a chain of thought style reasoning --- before generating the score for an acronym for a specific criteria, we generate a reasoning chain explicitly stating the reason for the scores. We use human evaluation to judge the final quality of the acronyms. An example of generated acronyms and associated feedback is given in \autoref{tab:acronym:feedback}.


\begin{table*}[h!]
\begin{adjustbox}{max width=\textwidth}
\small
\centering
\begin{tabular}{lp{4.8cm}p{5.8cm}}
\toprule
\textbf{Criteria} & output from GPT3: \textbf{STSLWN} & output from \ours: \textbf{Seq2Seq} \\ \midrule
Ease of pronunciation & Pronounced as ess-tee-ess-ell-double-you-enn which is very difficult. & Pronounced as seq-two-seq which is easy. \\ \midrule
Ease of spelling & Very difficult to spell. & Easy to spell. \\ \midrule
Relation to title & No relation to the title. & Mentions sequence which is somewhat related to the title.\\ \midrule
Positive connotation & Meaningless acronym. & Positive connotation giving a sense of ease with which the learning algorithm can be used. \\ \hline
Well-known & Not a well-known acronym. & Close to the word sequence which is a well-known word. \\ \midrule
Total score & 5/25 & 20/25 \\ \bottomrule
\end{tabular}
\end{adjustbox}
\caption{Comparison of acronyms for input = ``Sequence to Sequence Learning with Neural Networks''}
\label{tab:acronym:feedback}







\end{table*}
